\section{Noisy-Label \cifar{}}
\label{sec:results-cifar-noisy}

The noisy-label \cifar{} experiments test whether \method{} is useful when the training labels are corrupted but validation labels remain clean. The evaluation includes 20\% asymmetric label noise, 20\% symmetric label noise, and 40\% symmetric label noise. These settings are intended to create a mismatch between minimizing the corrupted training objective and improving validation accuracy on clean labels. Therefore, the distinction between peak validation accuracy, final validation accuracy, and best-to-final drop is central to the interpretation.

\subsection{Asymmetric 20\% Label Noise}

\input{tables/cifar_noisy_asym20}

\Cref{tab:cifar_asym20_main} shows that \method{} improves peak validation accuracy over the corresponding flat baseline for both \adamw{} and \sgdm{} under 20\% asymmetric label noise. With \adamw{}, peak accuracy increases from $58.4\%$ for the flat baseline to $61.7\%$ with \method{}, while Cosine+\method{} reaches $62.1\%$. With \sgdm{}, the flat baseline reaches $57.3\%$, while \method{} reaches $62.5\%$.

However, the final-accuracy results show that scheduler choice remains important. For \adamw{}, Cosine+\method{} gives the highest peak accuracy but finishes slightly below the cosine baseline, with final accuracies of $59.8\%$ and $60.0\%$, respectively. For \sgdm{}, cosine scheduling is the strongest configuration in both peak and final accuracy, reaching $63.3\%$ peak and $63.1\%$ final accuracy. Thus, under asymmetric label noise, \method{} improves the best checkpoint found by the flat optimizer, but does not uniformly improve final-checkpoint performance over conventional cosine scheduling.

\FloatBarrier
\subsection{Symmetric 20\% Label Noise}

\input{tables/cifar_noisy_sym20}

\Cref{tab:cifar_sym20_main} shows a similar pattern under 20\% symmetric label noise. With \adamw{}, \method{} improves peak validation accuracy from $58.6\%$ to $62.2\%$. When combined with cosine scheduling, \method{} reaches a similar peak of $62.0\%$ while improving final accuracy relative to \method{} alone, from $55.6\%$ to $58.4\%$. This suggests that cosine scheduling helps retain more of the performance found by adaptive learning-rate modulation.

For \sgdm{}, \method{} improves peak accuracy from $57.0\%$ to $62.5\%$, while Cosine+\method{} reaches the highest peak accuracy at $63.1\%$. However, cosine scheduling alone gives slightly higher final accuracy than Cosine+\method{}, with $62.0\%$ compared with $61.4\%$. This mirrors the asymmetric-noise setting: adaptive modulation improves the best checkpoint, while cosine scheduling remains important for retaining performance late in training.

\FloatBarrier
\subsection{Symmetric 40\% Label Noise}

\input{tables/cifar_noisy_sym40}

The 40\% symmetric label-noise setting makes the peak-versus-final distinction especially clear. As shown in \Cref{tab:cifar_sym40_main}, \method{} improves peak validation accuracy over the corresponding flat baseline for both optimizers. With \adamw{}, peak accuracy increases from $51.7\%$ to $55.0\%$. With \sgdm{}, it increases from $50.1\%$ to $55.6\%$.

At the same time, \method{} alone suffers substantial late-training degradation. For \adamw{}, \method{} reaches $55.0\%$ peak accuracy but falls to $39.9\%$ final accuracy, corresponding to a $15.1$ percentage-point drop. For \sgdm{}, \method{} reaches $55.6\%$ peak accuracy but falls to $39.7\%$ final accuracy, with a $16.0$ percentage-point drop. These results show that adaptive learning-rate modulation can find high-performing checkpoints under stronger label noise, but is not by itself sufficient to stabilize late training.

Cosine scheduling remains important for final-checkpoint stability. With \adamw{}, Cosine+\method{} improves peak accuracy over cosine alone, from $51.7\%$ to $54.7\%$, while giving similar final accuracy, $42.5\%$ compared with $42.2\%$. With \sgdm{}, Cosine+\method{} reaches the highest peak accuracy at $56.0\%$, but cosine alone gives slightly higher final accuracy, $46.9\%$ compared with $46.3\%$. This reinforces the interpretation that \method{} improves peak performance, whereas conventional decay schedules are still needed for stable final checkpoints.

\begin{figure}[t]
    \centering
    \thesisfigure{0.92\textwidth}{figures/cifar/cifar_symmetric40_curves.pdf}
    \caption[
        Validation trajectories on \cifar{} with 40\% symmetric label noise
    ]{
        Validation trajectories on \cifar{} with 40\% symmetric label noise. The figure illustrates the peak-versus-final behavior: \method{} reaches higher validation peaks than the corresponding flat baselines, but conventional decay schedules are important for late-training stability.
    }
    \label{fig:cifar-sym40-curves}
\end{figure}

\FloatBarrier
\subsection{Peak-Checkpoint Diagnostics}
\label{sec:cifar-noisy-checkpoint-diagnostics}

\input{tables/checkpointing_main}

To better understand what the peak checkpoints represent,
\Cref{tab:checkpoint-diagnostics-main} reports additional diagnostics at the
peak-validation checkpoint for the 20\% asymmetric and 40\% symmetric settings. These
two settings were chosen because they represent complementary corruption regimes:
asymmetric noise introduces structured, superclass-local corruptions, whereas 40\%
symmetric noise creates a severe random-corruption setting with pronounced
late-training degradation.

The diagnostics separate performance on corrupted training examples from performance
on uncorrupted training examples. For the corrupted subset, each example has two
relevant labels: the corrupted label used during training and the original clean label
before corruption. ``Corr. vs noisy'' measures accuracy on corrupted examples against
their corrupted labels, so high values indicate that the model predicts the imposed noisy
targets on the corrupted subset. ``Corr. vs clean'' evaluates the same corrupted examples
against their original clean labels, so higher values indicate that the model preserves
the underlying clean-label structure despite the corrupted training target. ``Clean vs
clean'' measures accuracy on the uncorrupted training examples and is included as a
sanity check that the model still learns the clean portion of the training set. Detailed
per-seed values are reported in \Cref{tab:checkpoint-diagnostics-detailed}.

The diagnostics show that the meaning of corrupted-subset accuracy depends strongly on
the noise process. Under 20\% asymmetric label noise, the cosine-only configurations
reach their peak-validation checkpoints much later. They also show near-complete
prediction of the corrupted labels: ``Corr. vs noisy'' is close to $100\%$, while
``Corr. vs clean'' is close to $0\%$. In contrast, \method{} and Cosine+\method{} reach
their peak-validation checkpoints much earlier and retain substantially higher accuracy
against the original clean labels on the corrupted subset. Together with the higher peak
validation accuracies, this suggests that the \method{} peak checkpoints are not merely
earlier checkpoints, but checkpoints that retain more clean-label structure while
achieving stronger validation performance.

Under 40\% symmetric label noise, the pattern is different. Because corrupted labels are
uniformly sampled from many incorrect classes, corrupted-target accuracy remains low
for all methods at the peak-validation checkpoints and is therefore a weaker direct
indicator of memorization than in the asymmetric setting. In this regime, ``Corr. vs
clean'' is more informative. \method{} improves peak validation accuracy over the flat
baseline and tends to improve clean-label accuracy on the corrupted subset, but
cosine-based configurations often provide stronger or comparable corrupted-subset
diagnostics while also giving better final-checkpoint stability. This supports the main
interpretation: \method{} helps find high-performing checkpoints, but conventional decay
schedules remain important for controlling late-training behavior.

The diagnostics in \Cref{tab:checkpoint-diagnostics-main} are computed at the
peak-validation checkpoints, because the purpose is to characterize the checkpoints that
produce $\accbest$. Final-epoch diagnostics are included in the detailed appendix table,
\Cref{tab:cifar_noisy_detailed_appendix}. These final diagnostics answer a different
question: they characterize the model after the full training horizon rather than at the
checkpoint selected by validation accuracy. Across many noisy-label runs, continued
training leads to much stronger prediction of the corrupted labels on the corrupted
subset, with ``Corr. vs noisy'' often approaching $100\%$ and ``Corr. vs clean'' often
approaching $0\%$. This is particularly clear for the cosine-based configurations, but is
not limited to them. The contrast supports the interpretation that peak-checkpoint
diagnostics and final-checkpoint diagnostics capture different stages of training: the
former describes the best checkpoint found during training, while the latter reflects the
extent to which continued training aligns the model with the corrupted training
objective.

\FloatBarrier
\subsection{Summary of Noisy \cifar{} Results}

Across the noisy \cifar{} experiments, \method{} consistently improves peak validation
accuracy over the corresponding flat baseline for both \adamw{} and \sgdm{}, and across
both symmetric and asymmetric label noise. This is the strongest evidence that adaptive
learning-rate modulation improves peak-checkpoint discovery in the tested vision
setting. The final-accuracy picture is weaker: \method{} on its own can suffer
substantial late-training degradation, and matching the best final-accuracy baselines
requires combining \method{} with a conventional decay schedule (see the subsections
above).

To test whether the peak gains can be explained by a better static learning-rate scale,
\Cref{tab:cifar_noisy_fixed_lr_ablation_summary} summarizes fixed learning-rate
multiplier ablations across the noisy \cifar{} settings. For each noise setting and
optimizer, the table reports the best fixed multiplier according to peak validation
accuracy, using the same multiplier candidates available to the \method{} controller.

\begin{table}[t]
\centering
\small
\setlength{\tabcolsep}{3.5pt}
\resizebox{\textwidth}{!}{%
\begin{tabular}{llccccc}
\toprule
\multirow{2}{*}{Noise} & \multirow{2}{*}{Optimizer} & \multirow{2}{*}{Best fixed} & \multicolumn{4}{c}{Validation best / final / drop (pp)} \\
& & & Best fixed & AEES & Cosine & Cosine + AEES \\
\midrule
Asym. 20\% & AdamW & $\times$1.0 & 58.4 / 56.2 / 2.2 & 61.7 / 58.0 / 3.7 & 60.2 / 60.0 / 0.2 & 62.1 / 59.8 / 2.3 \\
 & SGD & $\times$0.5 & 57.5 / 53.3 / 4.2 & 62.5 / 58.4 / 4.2 & 63.3 / 63.1 / 0.2 & 62.7 / 62.5 / 0.2 \\
\midrule
Sym. 20\% & AdamW & $\times$0.5 & 58.8 / 54.4 / 4.4 & 62.2 / 55.6 / 6.6 & 58.6 / 58.0 / 0.6 & 62.0 / 58.4 / 3.6 \\
 & SGD & $\times$0.5 & 57.4 / 50.6 / 6.8 & 62.5 / 56.4 / 6.1 & 62.2 / 62.0 / 0.2 & 63.1 / 61.4 / 1.6 \\
\midrule
Sym. 40\% & AdamW & $\times$0.5 & 52.5 / 39.0 / 13.5 & 55.0 / 39.9 / 15.1 & 51.7 / 42.2 / 9.5 & 54.7 / 42.5 / 12.2 \\
 & SGD & $\times$0.5 & 50.8 / 34.9 / 16.0 & 55.6 / 39.7 / 16.0 & 50.1 / 46.9 / 3.2 & 56.0 / 46.3 / 9.7 \\
\bottomrule
\end{tabular}
}
\caption{Fixed learning-rate multiplier ablation summary for noisy \cifar{}. For each noise setting and optimizer, the table reports the best fixed multiplier according to peak validation accuracy, using the same multiplier set available to the AEES controller. Validation entries are reported as best / final / drop, where drop is the best-to-final validation decrease in percentage points. This table tests whether the AEES peak gains under label noise exceed the best static LR rescaling baseline.}
\label{tab:cifar_noisy_fixed_lr_ablation_summary}
\end{table}

The fixed-multiplier results show that static rescaling alone does not explain the
\method{} peak gains. In every noisy setting in the table, \method{} reaches a higher
peak validation accuracy than the best fixed multiplier for the corresponding optimizer.
For example, under 40\% symmetric label noise, the best fixed multiplier reaches
$52.5\%$ peak accuracy with \adamw{} and $50.8\%$ with \sgdm{}, while \method{}
reaches $55.0\%$ and $55.6\%$, respectively. This supports the interpretation that
episodic adaptation contributes beyond selecting a single static learning-rate multiplier.

However, the gains are primarily peak-accuracy gains. \method{} alone often reaches its
best validation accuracy early and then degrades, especially under stronger symmetric
label noise. Final validation accuracy is more dependent on the presence and type of
global learning-rate schedule. Cosine scheduling often provides stronger final-checkpoint
stability, and in several \sgdm{} configurations it remains the best final-accuracy method.

The noisy \cifar{} results therefore support a bounded interpretation: \method{} is
effective as an adaptive mechanism for finding high-performing checkpoints under
corrupted supervision, and this effect is not explained by a single fixed learning-rate
multiplier. The checkpoint diagnostics further suggest that, especially under asymmetric
noise, these peaks are reached before full prediction of the corrupted targets and retain
more accuracy with respect to the original clean labels on the corrupted subset.
However, \method{} is not a standalone solution to noisy-label training and should not
be interpreted as a replacement for conventional decay schedules.